\chapter{Open Questions and Research Agenda}

\epigraph{\emph{The future of the program is not one horizon but several distances.}}{}

After the convergences of Part V and the demarcations of Chapter 14, the program no longer needs another chapter proving that it is allowed to ask its question. What it needs is an agenda. Harmonic Information Theory now has enough internal structure that its open problems can be ordered rather than merely accumulated. Some belong to the immediate continuation of active fronts. Some extend those fronts into new domains or new instruments. Others remain theoretical or technological horizons whose relevance is already visible even if their implementation is still distant. Writing the agenda clearly therefore matters for the same reason that writing the evidence clearly mattered earlier. It prevents the program from flattening near tasks, medium-range expansions, and speculative futures into one undifferentiated field of desire.

That ordering also protects the chapter from two opposite distortions. One would be to reduce the future of HIT to laboratory housekeeping: more seeds, more runs, more datasets, more measurements, as if the problem were simply to scale an already finished framework. The other would be to jump too quickly from convergence to grand application or metaphysics. The agenda has to hold both pressures at once. It has to say what should be done next, what new domains can now be entered, and what stronger formalizations or architectures become thinkable only because the earlier work has made them intelligible. In that sense, the chapter is neither a backlog nor a manifesto. It is the point where a distributed set of findings becomes a programmatic map.


\section{Why an agenda is needed now}


The most immediate reason for an explicit agenda is that HIT has entered a new phase of internal asymmetry. Some results are already strong and comparative. Others are real but still thinly seeded. Others belong to experimental practices whose protocolization is only now becoming sharp enough for public articulation. Without an ordered agenda, those unequal maturities are too easily misread. A frontier with a living computational probe, an acoustic-physiological probe, and a psychological field of practice can begin to look either more finished than it is or more chaotic than it is. The chapter must therefore preserve a simple discipline: every task should be located at the scale where it actually belongs.

This is also why the agenda must be written as a sequence and not as an inventory. The first horizon is immediate and programmatic: open questions that can already be attacked within active fronts. The second horizon is expansive: new domains, new carriers, and new descriptor families that grow directly out of the current work. The third horizon is formal and technological: stronger mathematical articulation, broader theoretical synthesis, and possible architectures that become imaginable only once the first two horizons have been clarified. That hierarchy does not make the later horizons less important. It makes them readable without inflation.


\section{Immediate questions and programmatic next steps}


Some of the next tasks are close enough to the existing program that they should be described not as future dreams but as direct continuations of work already under way. The first is comparability itself. Where a front still depends on single-seed or low-seed results, multi-seed validation remains a priority, not because repetition is bureaucratic virtue, but because the strongest claims of HIT depend on distinguishing structural signal from accidental optimization. The framework has already insisted that observation, hypothesis, and inference must not be collapsed into one another. The immediate agenda simply applies that same rule to active fronts: what has been seen strongly enough to support a conclusion, what still needs broader validation, and what remains suggestive but not yet settled.

The clearest immediate question in that sense is now \texttt{Speech$\leftrightarrow$EGG}. The point of the foundation-encoder regime test is not merely to swap one backbone for a larger one. It is to ask whether the legibility of descriptor-guided structure depends in part on the regime of representation itself. If a frozen foundation encoder on the speech side makes the descriptor signal newly visible, then the negative or null result of earlier arms must be reread as a question about representational power and asymmetry, not only about descriptors. If it does not, the null becomes harder and more informative. In both directions, the task is valuable because it reduces ambiguity left by the second front rather than multiplying novelty for its own sake.

An equally immediate task remains on the retrospective side of Escalon 1. The distinction between descriptor and mechanism was one of the major clarifications of the current manuscript, and it still has to be closed as an experimental reading rather than left as a conceptual correction. The remaining work there is not to reopen the entire archive. It is to consolidate the descriptor-by-mechanism contrast tightly enough that later generalizations to new domains are not built on a confusion the program itself has already diagnosed.


\section{Phideus beyond the current fronts}


Beyond those immediate continuations, the most important medium-range expansion is the Lissajous line. It matters because it is the first domain in which ratio can be generated, seen, and measured under exact control while also becoming shareable across the computational and analog branches of the program. On the Phideus side, that means retrieval between \texttt{audio XY} and figure, recovery of generative parameters, new topological and geometric descriptors, and a direct comparison between rational and irrational sweeps. On the Beacon side, it opens the possibility of a captured analog dataset built from physically excited media. The long-term importance of this front lies in the fact that it brings a problem usually scattered across acoustics, geometry, optics, and signal analysis into a single object with deterministic ground truth and cross-probe accessibility.

The same front also reopens an older and harder question: how much of a visible or measurable pattern belongs to the medium, how much to the excitation, and how much to the interaction between them. That is why the inverse-nodal problem matters here. The task is not only to classify figures or recover ratios. It is to ask how far one can infer the structure of the system from the topology of the pattern it stabilizes, and where that inference begins to fail under drift, nonlinearity, or irrational forcing \parencite{pinasco_2021,maynard_1985}. Read that way, Lissajous is not a decorative side branch. It is a controlled arena in which the relation between proportion, topology, inference, and medium can be pushed far harder than in the current audio-only benchmarks.

Phideus also has to widen what counts as a descriptor. The recurrence family already points in that direction, but recent work on multifrequency oscillatory portraits suggests an even broader route. Bahuguna and colleagues show that the interdependence structure of oscillatory portraits can become a predictive object in its own right rather than a mere summary of signals after the fact \parencite{bahuguna_2025}. For HIT, that result matters twice. It offers a template for descriptor families that operate trial by trial on relational organization across bands or oscillators, and it suggests a complementary metric space for evaluation beyond retrieval alone. A learned representation could then be asked not only whether it retrieves the correct partner, but whether it preserves or reorganizes the oscillatory portrait of the phenomenon under controlled transformations.

Two additional expansions follow from that same logic. One is transfer: whether a model trained in one cross-domain regime carries anything useful into another. The other is domain shift beyond acoustics. The cardiovascular line and the structural-vibration line are both important because they test whether the descriptor-guided logic survives when the carrier changes. Chakraborty and colleagues' structural vibration dataset is especially useful in that respect because it shows that informative identity can be embedded in oscillation propagating through built structures rather than through air alone \parencite{chakraborty_2025}. If harmonic information is really about structured relation and not about sound as such, then the move into vibration fields, sensor arrays, and technical media is not a decorative expansion. It is one of the strongest tests of the framework's stated generality.


\section{Beacon and PMP+Beacon: from observation to protocol}


The Beacon agenda begins from a different maturity profile. Here the problem is less the absence of interesting effects than the need to stabilize protocols around observations that are already rich. The corazofono line needs exactly that kind of transition. What matters is no longer only that a thoracic resonant center seems individually detectable under guided production, but how to measure its variation, compare personalized against generic settings, and establish what changes under exposure and calibration. The same is true of the EEG-to-harmonics line. Its importance lies not in presenting a spectacular real-time loop for its own sake, but in asking whether physiological streams can become control surfaces for harmonic synthesis in ways that remain legible, repeatable, and interpretable.

The pareidolia line belongs to the same family of problems. It does not need to be defended as if the experience itself were illusory until an instrument redeems it. Nor should it be inflated into a finished theory of perception. What the agenda requires is a disciplined protocol for distinguishing spontaneous projection, structurally induced projection, field dependence, and individual variation. The same discipline applies to the contrast between personalized and generic Beacon configurations. If the field matters differently when it is tuned to the organization of a particular body or voice, then that difference has to be rendered protocol-bearing rather than left as a persuasive anecdote.

The Lissajous convergence reappears here in its analog form. The construction of a captured analog dataset through physically excited media, optical readout, and systematic exploration is important not only because it links Beacon and Phideus. It also forces the Beacon line to formulate its own standards of control, drift management, and reproducibility. That makes it one of the strongest candidates for a shared experimental object across branches of the program.

The \texttt{PMP+Beacon} agenda begins from another asymmetry again. Here there is already repeated practice, interviews, before/after tests, and convergent observation; what remains underdeveloped is the comparative design capable of making those materials readable across cases. The central next step is therefore not to ask whether anything happens in sessions. It is to compare \texttt{PMP} alone with \texttt{PMP+Beacon} under conditions strong enough to say what changes in depth of access, sustained tension, symbolic elaboration, and later integration. That requires not only bio-signals and questionnaires, but also disciplined content analysis, judges blind to condition where possible, and a more explicit treatment of sample size and variation than the current practice alone can provide.

This is also where the language of registers matters again. A psychological field of practice cannot be reduced to physiology simply because physiology is measured alongside it, and symbolic transformation cannot be treated as noise because it resists easy quantification. The agenda is strongest when it keeps those registers distinct while still putting them into a shared protocol. That is the real task of \texttt{PMP+Beacon}: not to collapse therapeutic process into bio-signal, but to build a comparative architecture in which experiential, physiological, and symbolic change can be read together without being mistaken for the same type of evidence.


\section{Formalization and long-horizon technology}


The experimental agenda alone is not enough. HIT also needs stronger formalization of the concepts it has used to organize its evidence. The most obvious unfinished problem is harmonic efficiency itself. What kind of information is preserved, compressed, stabilized, or made easier to recover under simpler proportional relations? How much of that depends on the receiver, the medium, or the task? Questions of this sort connect the framework naturally to the free energy tradition, where the issue is not whether a system “likes” a field, but whether certain relations reduce corrective burden, prediction error, or coordination cost in ways that can be described without metaphysical inflation \parencite{friston_2010}. They also connect to biosemiotics, where the relation between sign, organism, and environment becomes relevant precisely when patterned differences begin to matter for coordination and action \parencite{hoffmeyer_2008,barbieri_2008}.

Another unfinished formal problem lies in the relation between structural coupling and mode locking. The analogy already plays an important role in the manuscript, but it still needs a more exact mathematical articulation. Can one formulate a metric of successful coupling that links the width of an Arnold tongue, the stability of a relational regime, and the properties of the receiving system without pretending that oscillators and organisms are the same kind of thing? The same question extends into the problem of autonomy and precariousness raised earlier in the manuscript: whether systems under different forms of fragility respond differently to the same harmonic field, and whether that difference can be made measurable without erasing the specificity of the organismic case \parencite{dipaolo_2005,maturana_varela_1987}.

On the side of descriptor design, Wang and colleagues' Harmonic Aggregation Entropy is useful precisely because it suggests a family of estimators that begin from consistency of relation rather than from isolated peaks or absolute magnitudes \parencite{wang_haagen_2025}. The point is not that Phideus should simply import that metric unchanged. It is that the space of possible descriptors is larger than interval summaries, spectral envelopes, or recurrence counts alone. Higher-order relational estimators, aggregation-based measures, and noisy-series harmonicity scores all point toward a richer future vocabulary for the computational branch.

The final horizon belongs to technology, but only after the experimental and formal questions have been differentiated clearly enough to support it. If Phideus continues to develop as a proportional sensor and Beacon as a proportional actuator, then the possibility of a more integrated architecture becomes newly legible. That is the point at which the name \texttt{PPU} can enter without sounding either premature or empty. A Proportional Processor Unity would not name a finished machine. It would name the horizon in which reading, comparing, and modulating proportional organization become coupled operations within one system. For now that remains speculative, and it should remain speculative. But it is no longer unintelligible. The agenda has made it sayable. The next chapter can therefore ask a different kind of question: not what must still be investigated, but what kinds of derivation, application, and public consequence might follow if this program continues to mature.
